\section{Introduction}
\label{sec:intro}

\subsection{Industrial Measurement and Control Context}

During load following and deep peak shaving, ultra-supercritical power units operate close to equipment protection limits. Their boiler-turbine coordinated control systems must keep pressure, separator enthalpy, and power within bounds while coordinating fuel, feedwater, and turbine-valve actuation. For this plant class, constraint handling is a coordinated control problem, not only a nominal tracking problem.

The controller design model is imperfect in service. Coal quality varies across shipments, heat absorption changes as heat-transfer surfaces foul, valve and fuel-delivery characteristics drift, and unmodeled coupling becomes more visible during load following. These effects create persistent mismatch between the controller model and the true process response. A controller that appears feasible under the nominal model, whether PID, LQR, MPC, or a learned policy, can still drive constrained variables toward a pressure, enthalpy, or power violation.

Recent applied-control studies similarly report disturbance response and computational update burden alongside tracking accuracy when plant response departs from the nominal design~\cite{wu2026pressure,feng2026event}. We adopt the same discipline for boiler-turbine model mismatch by treating safety-margin selection as a commissioning task for a downstream safety-filter overlay.

\subsection{Safety-Filter View}

We study a downstream safety filter, not a replacement for the plant controller. The filter receives current plant measurements and the upstream controller's proposed command, checks that command against process safety constraints, and returns the nearest admissible command. This placement suits industrial control practice because it separates tracking from protection: the existing controller handles load following and economic operation, while the safety layer enforces the final operating constraints.

We implement the projection as a low-latency quadratic program (QP) whose rows are derived from control barrier functions (CBFs)~\cite{ames2017cbf}. Relative degree must be defined with respect to the controller command rather than an unmeasured actuator state. After the fuel, feedwater, and turbine-valve dynamics are included, each pressure, separator-enthalpy, and power limit used here has command-level relative degree two. High-order CBFs (HOCBFs)~\cite{xiao2021hocbf} are therefore needed to preserve the constraint dynamics. A first-order CBF baseline would incorrectly assume that the commanded input appears in the first barrier derivative; that structural error is distinct from the plant-model mismatch studied below.

\subsection{Technical Gap: Uncertainty Propagation in High-Relative-Degree Constraints}

Gaussian processes (GPs) provide a natural way to learn model residuals from commissioning or operation data while quantifying posterior uncertainty~\cite{rasmussen2006gp,srinivas2010gp}. Existing GP-CBF methods usually apply the uncertainty bound to first-order barrier constraints~\cite{cheng2019guaranteed,fisac2018general,castaneda2021pointwise}, while recent GP-HOCBF work extends learning-based safety to high-order constraints through chance-constraint or convex reformulations~\cite{aali2024learning,zhang2025adaptive}. These studies establish useful theory, but they do not by themselves resolve the commissioning problem faced by actuator-limited boiler-turbine loops.

For boiler-turbine control, the difficulty is structural. The selected pressure, separator-enthalpy, and power limits are command-level relative-degree-two constraints, and residual uncertainty propagates through the HOCBF recursion rather than appearing only at the final constraint row. Applying the full implemented margin can over-tighten the online QP, reduce actuator authority, or trigger fallback behavior. Applying no margin can also be unsafe when the perturbation is state-dependent and amplifies along the trajectory. The unresolved question is therefore not only how to construct a GP-HOCBF margin, but how much of that margin should be applied in real time for the residual structure observed during commissioning.

RoCBF-SF addresses this question by combining GP residual mean correction, compositional uncertainty propagation through the HOCBF $\psi$-chain, QP projection, and a robustness factor calibrated from commissioning residuals. The scientific contribution is the safety-filter overlay and its calibration procedure under measured model mismatch, rather than a new upstream controller.

\subsection{Contributions}

We introduce RoCBF-SF, a commissioning-calibrated GP-HOCBF safety filter for ultra-supercritical boiler-turbine control. The paper makes four contributions to industrial measurement, constrained control, and real-time deployment:

\begin{enumerate}
    \item \textbf{Downstream safety-filter architecture for boiler-turbine process control.} RoCBF-SF preserves the upstream controller while enforcing command-level relative-degree-two pressure, separator-enthalpy, and power constraints obtained from an actuator-augmented process model.

    \item \textbf{Compositional uncertainty propagation through the HOCBF chain.} The method propagates GP posterior residual uncertainty through each HOCBF $\psi$-level and embeds the resulting state-dependent margin in a real-time QP projection.

    \item \textbf{Commissioning-calibrated robustness factor.} The scalar factor $\epsilon_\kappa$ interpolates between mean-corrected operation and the full implemented uncertainty margin, with diagnostics for selecting this factor under additive and state-dependent mismatch.

    \item \textbf{Simulation validation and measured-operation evidence.} The approach is evaluated on an actuator-augmented boiler-turbine benchmark, while bounded production historian records and controller-log exports from a 660\,MW ultra-supercritical unit provide response, execution, feasibility, and latency evidence under measured operation.
\end{enumerate}

The validation is organized around this applied process-control logic. Controlled simulations first isolate the model-mismatch mechanism and show when GP mean correction and $\epsilon_\kappa$ calibration are needed. Plant historian records then compare pre- and post-retrofit response under matched load movement, and controller-log exports verify the online execution path across repeated post-retrofit windows. Together, these evidence sources connect the formal safety-filter mechanism to plant measurements while keeping the upstream controller outside the proposed method.
